% Lösungen Einheit 2 von 3 – Achsensymmetrie und Spiegeln (zusammenbau v0.1)
\einheitenkopf{Einheit 2 von 3 · Achsensymmetrie und Spiegeln}

\erg{16}{a) $2$ Achsen: die beiden Mittellinien}

16a)

\rechteck[achsen=beide]{5}{2}

%% TODO Lösung der Erklärzeile 17g) fehlt in der Bank
\erg{17}{a) ja \quad b) die senkrechte Diagonale durch die obere und die untere Spitze \quad c) die senkrechte Gerade von der Spitze zur Mitte der Grundseite \quad d) die Gerade durch die Mitten der beiden parallelen Seiten \quad e) die senkrechte Diagonale durch die obere und die untere Spitze \quad f) die senkrechte Gerade von der Spitze zur Mitte der Grundseite \quad g) \ldots \quad h) $1$ (Drachenviereck) \quad i) $2$ (die beiden Diagonalen) \quad j) Parallelogramm \quad k) $4$, $2$, $2$, $1$ (Quadrat, Rechteck, Raute, Drachenviereck) \quad l) $3$ \quad m) $2$ (senkrecht und waagerecht) \quad n) $AF = 6 : 2 = 3$\,cm \quad o) 1 \quad p) 1}

17b)

\drachen[achsen=senkrecht]{4}{3}{0.3}


17e)

\drachen[achsen=senkrecht]{5}{3}{0.25}

%% TODO Lösung der Erklärzeile 18g) fehlt in der Bank
\erg{18}{a) ja \quad b) $P'(6|3)$ \quad c) $P'(5|5)$ \quad d) $P'(2|4)$ \quad e) $P'(7|6)$ \quad f) $P'(7|1)$ \quad g) \ldots \quad h) Bilddreieck $A'(7|1)$, $B'(5|1)$, $C'(6|4)$ \quad i) Bilddreieck $A'(1|7)$, $B'(4|7)$, $C'(4|5)$ \quad j) $A'(-2|1)$, $B'(-5|1)$, $C'(-4|4)$ \quad k) ergänzte Ecken $B'(6|1)$ und $C'(6|4)$; es entsteht ein Fünfeck in Hausform \quad l) Bilddreieck $A'(3|1)$, $B'(5|2)$, $C'(5|1)$ \quad m) ein gleichschenkliges Dreieck \quad n) Drachenviereck (a ist eine Diagonale, die andere Diagonale steht senkrecht darauf und wird halbiert; keine Raute, weil das Dreieck nicht gleichschenklig ist)}
\erg{19}{a) nein}
\erg{20}{a) Quadrat, Rechteck, Raute}
\erg{21}{a) Nur $1$ Achse: die waagerechte passt nicht, weil die beiden parallelen Seiten verschieden lang sind. \quad b) Faltet man es an einer Diagonale oder einer Mittellinie, passen die Hälften nicht aufeinander – die Ecken liegen versetzt. \quad c) z. B. ein Drachenviereck (Achse: eine Diagonale) oder ein gleichschenkliges Trapez \quad d) Rechteck}
